\documentclass[11pt,a4paper,twoside]{article}
\usepackage[utf8]{inputenc}
\usepackage[english]{babel}
\usepackage{actuarialangle, actuarialsymbol}

\begin{document}
Define the notations and central objects for learning theory

The labeled domain space \( \mathcal{Z} := \mathcal{X} \times \mathcal{Y} \)
consists of the domain set \( \mathcal{X} \) consisting of features (coordinates)
and the label set \( \mathcal{Y} \), for purely theoretical considerations \( \mathcal{Y} = \{0, 1\} \).
The training data \( S = \{ (x_0, y_0), \dots, (x_n, y_n)\}_{n \in \mathbb{N}} \) is a finite sequence of domain points.
An algorithm (learner) outputs a prediction rule/predictor/hypothesis \( h : \mathcal{X} \to \mathbb{Y} \), so \( A(s) = h \)
from a set of possible hypotheses \( \mathcal{H} \).
There is some labelling function \( f : \mathcal{X} \to \mathcal{Y} \), which we'll try to predict.
If \( f \) is deterministic then we'll consider a distribution \( \mathcal{D} \) over \( \mathcal{X} \),
otherwise over \( \mathcal{Z} \) (in that way, we also try to estimate the distribution over the outputs
as they can randomly change). We'll typically assume that samples \( (x, y) \sim \mathcal{D} \) are i.i.d.

The error of the classifier or true error/risk is

\( L_{\mathcal{D},f}(h) = \mathbb{P}_{x \sim \mathcal{D}}[h(x) \neq f(x)] = \mathcal{D}(\{x : h(x) \neq f(x))\}\)
or if \( \mathcal{D} \) is over \( \mathcal{Z} \),
\( L_{\mathcal{D}}(h) = \mathbb{P}_{(x,y) \sim \mathcal{D}}[h(x) \neq y] = \mathcal{D}(\{x : h(x) \neq y)\}\).

This stands contrary to the empirical error/risk

\( L_S(h) = \frac{|i \in [m] : h(x_i) \neq y_i|}{m} \)
or \( f(x_i) \) instead of \( y_i \). 
Alternatively more general loss functions are possible.

The probability of a non-representative sample is given by \( \delta \)
viz a vis \( 1 - \delta \) the confidence parameter, this is simply \( \mathcal{D}(A), A \subset \mathcal{Z} \)
The accuracy parameter \( \varepsilon \) is typically used in statements of the form \( L_{\mathcal{D},f}(h) \leq \varepsilon \).



Describe the issue of overfitting and relate it to finite hypothesis classes

Assume a uniform distribution \( \mathcal{D} \) over \( [0, 2]^2 \) and labeling \( f(x) = 1 \) if \( x \) is in the 
inner square \( [\frac{1}{2}, \frac{3}{2}]^2 \), otherwise \( 0 \). For the picture consider the predictor

\[
h_S(x) = \begin{cases}
  y_i & \exists i \in [m], s.t. x_i = x \\
  0
\end{cases}
\]
then \( L_S(h) = 0 \), but of course \( L_{\mathcal{D},f}(h) \) is rather large.

This is overfitting and can be migitated by using a subclass \( \mathcal{H} \) of hypotheses of all possible. 
Thus instead of considering just the Empirical Risk Minimization ERM
\( \argmin_h L_S(h) \), we use ERM with inductive bias \( \text{ERM}_{\mathcal{H}} = \argmin_{h \in \mathcal{H}} L_S(h) \).


Define PAC-learnability and state the sample complexity bound for finite hypotheses classes.

Probably Approximately Correct (PAC) learning.
A hypothesis class \( \mathcal{H} \) is PAC learnable
if there exist a function \( m_{\mathcal{H}} : (0, 1)^2 \to \mathbb{N} \) and a learning algorithm with the
following property: For every \(\varepsilon, δ ∈ (0, 1)\), for every distribution \( \mathcal{D} \) over \( \mathcal{X} \) 
and for every labeling function \( f : \mathcal{X} \to \{0, 1\}\), if the realizable assumption holds
with respect to \( \mathcal{H}, \mathcal{D},f \), then when running the learning algorithm on \( m \geq m_{\mathcal{H}}(\varepsilon, \delta) \)
i.i.d. examples generated by \( \mathcal{D} \) and labeled by \( f \), the algorithm returns
a hypothesis \( h \) such that, with probability of at least \( 1 − δ \) (over the choice of the examples), 
\( L_{\mathcal{D},f} \leq \varepsilon \).

\( m_{\mathcal{H}}(\varepsilon, \delta) \) is also called the sample complexity.
The Realizability Assumption is the existance of some \( h^\ast \in \mathcal{H}, \text{s.t.} L_{(\mathcal{D},f)}(h^\ast) = 0 \)

Every finite \( \mathcal{H} \) is PAC learnable with 
\( m_{\mathcal{H}}(\varepsilon, \delta) \leq \text{ceil} \frac{\log(|\mathcal{H}|/\delta)}{\varepsilon} \)



State the Bayes optimal predictor, agnostic learnability and generalized loss functions

Given any \( \mathcal{D} \) over \( \mathcal{Z} \) it can be shown that
\( h_{\mathcal{D}}(x) = \begin{cases}
  1 & \mathbb{P}[y = 1|x] \geq \frac{1}{2} \\
  0
\end{cases} \)
is optimal in the sense that, \( L_{\mathcal{D}}(h_{\mathcal{D}} \leq L_{\mathcal{D}}(g) \) for all \( g \).

Agnostic PAC Learnability
is PAC Learnability but dropping the realizable assumption and instead asking for
\( L_{\mathcal{D}}(h) \leq \min_{h' \in \mathcal{H}} L_{\mathcal{D}}(h') + \varepsilon \)
with probability \( 1 - \delta \) and \( h \) resulting from training on at least \( m \geq m_{\mathcal{H}}(\varepsilon, \delta)\) samples.

The generalized loss function is
\( L_{\mathcal{D}}(h) = \mathbb{E}_{x \sim \mathcal{D}}[\el(h, z)] \),
where \( \el : \mathcal{H} \times \mathcal{Z} \to \mathbb{R}_{+} \) and the sample loss function
\( L_S(h) = \frac{1}{m} \sum_{i=1}^m \el(h, z_i) \).

What are ϵ-representative samples and how does it imply PAC-learnability?

(ϵ-representative sample) A training set \(S\) is called ϵ-representative (w.r.t. domain \(Z\), 
hypothesis class \(\mathcal{H}\), loss function \(\ell\), and distribution \(\mathcal{D}\)) if 
\(\forall h \in \mathcal{H}, \, |L_{S}(h) - L_{\mathcal{D}}(h)| \leq ϵ.\)

Assume that a training set \(S\) is \(\frac{ϵ}{2}\)-representative (w.r.t. domain \(Z\), 
hypothesis class \(\mathcal{H}\), loss function \(\ell\), and distribution \(\mathcal{D}\)). 
Then, any output of \(\text{ERM}_{\mathcal{H}}(S)\), namely, any \(h_{S} \in \arg\min_{h \in \mathcal{H}} L_{S}(h)\), 
satisfies 
\(L_{\mathcal{D}}(h_{S}) \leq \min_{h \in \mathcal{H}} L_{\mathcal{D}}(h) + ϵ.\)
Proof:
For every \( h \in \mathcal{H} \)
\( L_{\mathcal{D}}(h_S) \leq L_S(h_S) + \frac{\varepsilon}{2} \leq L_S(h) + \frac{\varepsilon}{2} \leq L_{\mathcal{D}}(h) + \varepsilon \).
where the first and third inequalities are due to the assumption that \( S \) is \( \frac{\varepsilon}{2} \) -
representative (Definition 4.1) and the second inequality holds since \( h_S \) is an ERM predictor.



Define uniform convergence and how is it related to PAC learnability

We say that \( \mathcal{H} \) has the uniform convergence property, if there exists 
\( m_{\mathcal{H}}^{UC} : (0, 1)^2 \to \mathbb{N} \) s.t. for every \( \varepsilon, \delta (0, 1) \)
and \( \mathcal{D} \) over \( \mathcal{Z} \), if \( S \) is an i.i.d. sample of size \( m \geq m_{\mathcal{H}}^{UC}(\varepsilon, \delta) \)
then with probability \( 1 - \delta \), \( S \) is a \( \varepsilon \)-representative.

If a class \(\mathcal{H}\) has the uniform convergence property with a function \(m_{\mathcal{H}}^{\text{UC}}\) 
then the class is agnostically PAC learnable with the sample complexity \(m_{\mathcal{H}}(ϵ, δ) \leq m_{\mathcal{H}}^{\text{UC}}(ϵ/2, δ)\). 
Furthermore, in that case, the \(\text{ERM}_{\mathcal{H}}\) paradigm is a successful agnostic PAC learner for \(\mathcal{H}\).

A useful fact in the proof of the next result is
(Hoeffding's Inequality) 
Let \(\theta_1, \ldots, \theta_m\) be a sequence of i.i.d. random variables and assume that for all \(i\), 
\(\mathbb{E}[\theta_i] = \mu\) and \(\mathbb{P}[a \leq \theta_i \leq b] = 1\). Then, for any \(ϵ > 0\)
\(
\mathbb{P} \left[ \left| \frac{1}{m} \sum_{i=1}^{m} \theta_i - \mu \right| > ϵ \right] \leq 2 \exp \left( -2mϵ^2 / (b-a)^2 \right).
\)

(Finite \(\mathcal{H} \implies \) uniform convergence property)
Let \(\mathcal{H}\) be a finite hypothesis class, let \(Z\) be a domain, and let \(\ell : \mathcal{H} \times Z \to [0,1]\) be a loss function. 
Then, \(\mathcal{H}\) enjoys the uniform convergence property with sample complexity
\(
m_{\mathcal{H}}^{\text{UC}}(ϵ, δ) \leq \left\lceil \frac{\log(2|\mathcal{H}|/δ)}{2ϵ^2} \right\rceil.
\)
Furthermore, the class is agnostically PAC learnable using the ERM algorithm with sample complexity
\(
m_{\mathcal{H}}(ϵ, δ) \leq m_{\mathcal{H}}^{\text{UC}}(ϵ/2, δ) \leq \left\lceil \frac{2\log(2|\mathcal{H}|/δ)}{ϵ^2} \right\rceil.
\)

State the no-free-lunch theorem for learnability and it's Corollary

(No-Free-Lunch) 
Let \(A\) be any learning algorithm for the task of binary classification with respect to the \(0 - 1\) loss over a domain \(\mathcal{X}\). 
Let \(m\) be any number smaller than \(|\mathcal{X}|/2\), representing a training set size. 
Then, there exists a distribution \(\mathcal{D}\) over \(\mathcal{X} \times \{0,1\}\) such that:
1. There exists a function \(f : \mathcal{X} \to \{0,1\}\) with \(L_{\mathcal{D}}(f) = 0\).
2. With probability of at least \(1/7\) over the choice of \(S \sim \mathcal{D}^m\), we have that \(L_{\mathcal{D}}(A(S)) \geq 1/8\).

Let \( C \) be a subset of \( \mathcal{X} \) of size \( 2m \). 
The intuition of the proof is that any learning algorithm that observes 
only half of the instances in \( C \) has no information on what should be the labels of the rest of the instances in \( C \).

Let \( \{f_1, \dots, f_T | f_i : C \to \{0, 1\} \},\ T = 2^{2m} \) and let \( D_i \) be a distribution over
\(C \times \{0,1\}\) defined by
\[
\mathcal{D}_i(\{(x,y)\}) = 
\begin{cases} 
1/|C| & \text{if } y = f_i(x) \\
0 & \text{otherwise}.
\end{cases}
\]
That is, the probability to choose a pair \((x, y)\) is \(1/|C|\) if the label \(y\) is indeed the true label according to \(f_i\), 
and the probability is \(0\) if \(y \neq f_i(x)\). Clearly, \(L_{\mathcal{D}_i}(f_i) = 0\).
It then can be shown that \( \text{max}_{i \in [T]} \mathbb{E}_{S \sim \mathcal{D}^m_i}[L_{\mathcal{D}_i}(A(S))] \leq \frac{1}{4}\)
(which implies with some effort that this is not PAC learnable), an adversary can thus construct always a suitable \( f \) which is optimal
but the learning algorithm fails.


Let \( \mathcal{X} \) be an infinite domain set and let \( \mathcal{H} \) be the set of all
functions from \( \mathcal{X} \) to \( \{0, 1\} \). Then, \( \mathcal{H} \) is not PAC learnable.


How can infinite-size classes be learned?

Let \( \mathcal{H} := \{ h_{a} := \mathbb{1}_{x < a} : a \in \mathbb{R} \} \), clearly \( |\mathcal{H}| = 2^\aleph \)
Still \( \mathcal{H} \) is PAC learnable, using the ERM rule, with sample complexity of \( m_{\mathcal{H}}(\varepsilon, \delta) \leq \text{ceil}(\log(\frac{2}{\delta}/\varepsilon)) \).

Let \( a^\ast \) be s.t. \( L_{\mathcal{D}}(h_{a^\ast}) = 0 \) and \( (a_0, a_1) \) s.t. 
\( \mathbb{P}_{x \sim \mathcal{D}_x}[x \in (a_0, a^\ast)] = \mathbb{P}_{x \sim \mathcal{D}_x}[x \in (a^\ast, a_1)] = \varepsilon \)
Given a training set \( S \), set \( b_0 = \max \{x : (x, 1) \in S\} \) and \( b_1 = \min\{x : (x, 0) \in S \}\).
Set \( h_{b_S} = \text{ERM}_S \), thus \( b_S \in (b_0, b_1) \) and 
\( \mathbb{P}_{S \sim \mathcal{D}^m}[L_{\mathcal{D}}(h_S) > \varepsilon] \leq \mathbb{P}_{S \sim \mathcal{D}^m}[b_0 < a_0 \vee b_1 > a_1] \)
would imply \( L_{\mathcal{D}}(h_{S}) \leq \varepsilon \). Note that
\( \mathbb{P}_{S \sim \mathcal{D}^m}[b_0 < a_0] = \mathbb{P}_{S \sim \mathcal{D}^m}[(x, y) \in S, x \nin (a_0, a^\ast)] = (1-\varepsilon)^m \leq e^{-\varepsilon m} \)
by assumption \( m > \log(\frac{2}{\delta})/\varepsilon  \), thus the bound simplifies to \( \frac{\delta}{2} \) the same can be done for \( b_1 > a_1 \) and the result follows.


Define shattering, give an example and explain why it prohibits PAC learnability

(Restriction of \(\mathcal{H}\) to \(C\)) 
Let \(\mathcal{H}\) be a class of functions from \(\mathcal{X}\) to \(\{0,1\}\) and let \(C = \{c_1, \ldots, c_m\} \subset \mathcal{X}\). 
The restriction of \(\mathcal{H}\) to \(C\) is the set of functions from \(C\) to \(\{0,1\}\) that can be derived from \(\mathcal{H}\). That is,
\(
\mathcal{H}_C = \{(h(c_1), \ldots, h(c_m)) : h \in \mathcal{H}\},
\)
where we represent each function from \(C\) to \(\{0,1\}\) as a vector in \(\{0, 1\}^{|C|}\).

(Shattering) A hypothesis class \(\mathcal{H}\) shatters a finite set \(C \subset \mathcal{X}\), 
if the restriction of \(\mathcal{H}\) to \(C\) is the set of all functions from \(C\) to \(\{0,1\}\). 
That is, \(|\mathcal{H}_C| = 2^{|C|}\).

Example:
Let \(\mathcal{H}\) be the class of threshold functions over \(\mathbb{R}\). Take a set \(C = \{c_1\}\). 
Now, if we take \(a = c_1 + 1\), then we have \(h_a(c_1) = 1\), and if we take \(a = c_1 - 1\), then we have \(h_a(c_1) = 0\). 
Therefore, \(\mathcal{H}_C\) is the set of all functions from \(C\) to \(\{0,1\}\), and \(\mathcal{H}\) shatters \(C\). 
Now take a set \(C = \{c_1, c_2\}\), where \(c_1 \leq c_2\). 
No \(h \in \mathcal{H}\) can account for the labeling \((0,1)\), because any threshold that assigns the label \(0\) to \(c_1\) 
must assign the label \(0\) to \(c_2\) as well. 
Therefore not all functions from \(C\) to \(\{0,1\}\) are included in \(\mathcal{H}_C\); hence \(C\) is not shattered by \(\mathcal{H}\).

In the proof of the free-lunch theorem was an adversary, who could construct a distribution which is realizable but not learnable,
if the domain could be shattered.


Define the VC-dimension and the fundamental theorem of PAC learning

The VC-dimension (Vapnik-Chervonenkis) of a hypothesis class \( \mathcal{H} \), denoted \( \text{VCdim}(\mathcal{H}) \), 
is the maximal size of a set \( C \subset X \) that can be shattered by \( \mathcal{H} \). 
If \( \mathcal{H} \) can shatter sets of arbitrarily large size we say that \( \mathcal{H} \) has infinite
VC-dimension.

It follows by the no-free-lunch theorem, that if \( \mathcal{H} \) has infinite VC-dimension. Then, \( \mathcal{H} \) is not PAC
learnable, since for any training set size \( m \), there is a shattered set of size \( 2m \).

To determine the VC-dimension, the steps are
1. There exists a set \( C \) of size \( d \) that is shattered by \( \mathcal{H} \).
2. Every set \( C \) of size \( d+1 \) is not shattered by \( \mathcal{H} \).

(The Fundamental Theorem of Statistical Learning) 
Let \(\mathcal{H}\) be a hypothesis class of functions from a domain \(\mathcal{X}\) to \(\{0,1\}\) and 
let the loss function be the \(0 - 1\) loss. Then, the following are equivalent:
1. \(\mathcal{H}\) has the uniform convergence property.
2. Any ERM rule is a successful agnostic PAC learner for \(\mathcal{H}\).
3. \(\mathcal{H}\) is agnostic PAC learnable.
4. \(\mathcal{H}\) is PAC learnable.
5. Any ERM rule is a successful PAC learner for \(\mathcal{H}\).
6. \(\mathcal{H}\) has a finite VC-dimension.

The tricky part is \( 6 \implies 1 \), the proof has the following 2 central ideas

(Growth Function) 
Let \(\mathcal{H}\) be a hypothesis class. 
Then the growth function of \(\mathcal{H}\), denoted \(\tau_{\mathcal{H}} : \mathbb{N} \to \mathbb{N}\), is defined as 
\(\tau_{\mathcal{H}}(m) = \max_{C \subset \mathcal{X}: |C|=m} |\mathcal{H}_C|.\)
In words, \(\tau_{\mathcal{H}}(m)\) is the number of different functions from a set \(C\) of size \(m\) to \(\{0,1\}\) 
that can be obtained by restricting \(\mathcal{H}\) to \(C\).

Sauer-Shelah-Perle
Let \( \mathcal{H} \) be a hypothesis class with \( \text{VCdim}(\mathcal{H}) \leq d < \infty \). 
Then, for all \( m \), \(\tau_{\mathcal{H}}(m) \leq \sum_{i=0}^d \binom{m}{i}\) . 
In particular, if \( m > d + 1 \) then \( \tau_{\mathcal{H}}(m) \leq (\frac{em}{d})^d \).


What is the error decomposition?

\(L_{\mathcal{D}}(h_S) = \epsilon_{\text{app}} + \epsilon_{\text{est}}\)
where: \(\epsilon_{\text{app}} = \min_{h \in \mathcal{H}} L_{\mathcal{D}}(h)\) 
\(\epsilon_{\text{est}} = L_{\mathcal{D}}(h_S) - \epsilon_{\text{app}}.\)

Approximation Error
This term measures how much risk we have because
we restrict ourselves to a specific class, namely, how much inductive bias we
have.
Estimation Error
The estimation error results
because the empirical risk (i.e., training error) is only an estimate of the
true risk.
\end{document}

